\documentclass[10pt,a4paper]{amsart}
\usepackage{amsmath,amssymb,amsthm,graphicx}
\usepackage{verbatim}
\newenvironment{reference}{\comment}{\endcomment}
%\newenvironment{reference}{}{}
\newenvironment{slogan}{\comment}{\endcomment}
\newenvironment{history}{\comment}{\endcomment}

% For commutative diagrams we use Xy-pic
\usepackage[all]{xy}

% We use 2cell for 2-commutative diagrams.
\xyoption{2cell}
% Theorem environments.
%
\theoremstyle{plain}
\newtheorem{theorem}[subsection]{Theorem}
\newtheorem{proposition}[subsection]{Proposition}
\newtheorem{lemma}[subsection]{Lemma}

\theoremstyle{definition}
\newtheorem{definition}[subsection]{Definition}
\newtheorem{example}[subsection]{Example}
\newtheorem{exercise}[subsection]{Exercise}
\newtheorem{situation}[subsection]{Situation}

\theoremstyle{remark}
\newtheorem{remark}[subsection]{Remark}
\newtheorem{remarks}[subsection]{Remarks}

\numberwithin{equation}{subsection}

% Macros
%
\def\lim{\mathop{\mathrm{lim}}\nolimits}
\def\colim{\mathop{\mathrm{colim}}\nolimits}
\def\Spec{\mathop{\mathrm{Spec}}}
\def\Hom{\mathop{\mathrm{Hom}}\nolimits}
\def\Ext{\mathop{\mathrm{Ext}}\nolimits}
\def\SheafHom{\mathop{\mathcal{H}\!\mathit{om}}\nolimits}
\def\SheafExt{\mathop{\mathcal{E}\!\mathit{xt}}\nolimits}
\def\Sch{\mathit{Sch}}
\def\Mor{\mathop{\mathrm{Mor}}\nolimits}
\def\Ob{\mathop{\mathrm{Ob}}\nolimits}
\def\Sh{\mathop{\mathit{Sh}}\nolimits}
\def\NL{\mathop{N\!L}\nolimits}
\def\CH{\mathop{\mathrm{CH}}\nolimits}
\def\proetale{{pro\text{-}\acute{e}tale}}
\def\etale{{\acute{e}tale}}
\def\QCoh{\mathit{QCoh}}
\def\Ker{\mathop{\mathrm{Ker}}}
\def\Im{\mathop{\mathrm{Im}}}
\def\Coker{\mathop{\mathrm{Coker}}}
\def\Coim{\mathop{\mathrm{Coim}}}

% Boxtimes
%
\DeclareMathSymbol{\boxtimes}{\mathbin}{AMSa}{"02}

%
% Macros for moduli stacks/spaces
%
\def\QCohstack{\mathcal{QC}\!\mathit{oh}}
\def\Cohstack{\mathcal{C}\!\mathit{oh}}
\def\Spacesstack{\mathcal{S}\!\mathit{paces}}
\def\Quotfunctor{\mathrm{Quot}}
\def\Hilbfunctor{\mathrm{Hilb}}
\def\Curvesstack{\mathcal{C}\!\mathit{urves}}
\def\Polarizedstack{\mathcal{P}\!\mathit{olarized}}
\def\Complexesstack{\mathcal{C}\!\mathit{omplexes}}
% \Pic is the operator that assigns to X its picard group, usage \Pic(X)
% \Picardstack_{X/B} denotes the Picard stack of X over B
% \Picardfunctor_{X/B} denotes the Picard functor of X over B
\def\Pic{\mathop{\mathrm{Pic}}\nolimits}
\def\Picardstack{\mathcal{P}\!\mathit{ic}}
\def\Picardfunctor{\mathrm{Pic}}
\def\Deformationcategory{\mathcal{D}\!\mathit{ef}}
\usepackage{cleveref}
\crefname{originalreference}{original reference}{original references}
\Crefname{originalreference}{Original reference}{Original references}
\makeatletter
\newlabel{lemma-divide-pullback-good-blowing-up}{{0FIH}{source}{}{sourcepointer-1}{}}
\newlabel{lemma-divide-pullback-good-blowing-up@cref}{{[originalreference][0][]0FIH}{[1][0][]source}{}{}{}}
\makeatother
\crefname{originalreference}{original reference}{original references}
\Crefname{originalreference}{Original reference}{Original references}
\makeatletter
\newlabel{lemma-enough-better}{{0FIR}{source}{}{sourcepointer-2}{}}
\newlabel{lemma-enough-better@cref}{{[originalreference][0][]0FIR}{[1][0][]source}{}{}{}}
\makeatother
\crefname{originalreference}{original reference}{original references}
\Crefname{originalreference}{Original reference}{Original references}
\makeatletter
\newlabel{lemma-ok-for-projective-bundle}{{0FIN}{source}{}{sourcepointer-3}{}}
\newlabel{lemma-ok-for-projective-bundle@cref}{{[originalreference][0][]0FIN}{[1][0][]source}{}{}{}}
\makeatother
\crefname{originalreference}{original reference}{original references}
\Crefname{originalreference}{Original reference}{Original references}
\makeatletter
\newlabel{lemma-pushforward-blowup}{{0FIS}{source}{}{sourcepointer-4}{}}
\newlabel{lemma-pushforward-blowup@cref}{{[originalreference][0][]0FIS}{[1][0][]source}{}{}{}}
\makeatother
\title{Compatibility of the cycle class map with pushforwards}
\author{The Stacks Project Authors}
\date{stacks-project:weil.tex:local-snapshot-20261012; GNU FDL 1.2 or later}
\begin{document}
\maketitle
\noindent\textit{Editorial note.} All delivered mathematical text is the author's own, copied from the pinned Stacks Project chapter file weil.tex; no statement or proof step was rewritten, reordered or removed.\par\medskip
\noindent\textit{Editorial note.} The counted proof is the author's complete argument: the reduction through the blowing up, the use of the compatibility for projective space bundles and of axiom (A9), and the final computation through the projection formula.\par\medskip
\noindent\textit{Editorial note.} The label lemma-done names the counted statement, and the proof environment immediately following it in the source is delivered verbatim as its proof.\par\medskip
\noindent\textit{Editorial note.} The counted statement is carried by the author's own lemma environment in the chapter source.\par\medskip
\noindent\textit{Editorial note.} Pre-open gate: python operations/agent-authoring/tools/precard\_gate.py --source weil.tex --label lemma-done --start 5384 --end 5419 --local-id 28153 returned RESULT=OK (L28153 free, occupancy FREE, public interval clear). Pre-assembly tag self-check: canonical tag resolved and not present in reference/public-items.jsonl.\par\medskip
\noindent\textit{Editorial note.} The statement and proof are copied from the pinned chapter file; the cited inputs are kept as external pointers with their canonical Stacks tags.\par\medskip
\noindent\textit{Allowed standard prerequisites (editorial).}\begin{itemize}
\item The compatibility for projective space bundles. Same chapter, the lemma that the cycle class map commutes with pushforward along the projectivization of a vector bundle. Supplies the equality expressing the pushforward of the cycle class of a pulled back cycle as the unit.
\item Compatibility of the cycle class map with the class of a divisor. Axiom (A9) describing the cycle class of the fundamental class of an effective Cartier divisor as the pushforward of the unit. Identifies the pushforward of the unit along the closed immersion with the cycle class of the divisor.
\item The pushforward of the unit along the blowing up is one. Same chapter, the lemma computing the pushforward of the unit along the blowing up appearing in the reduction. Shows that the pushforward of the unit along the morphism is one, which finishes the computation.
\end{itemize}\medskip
\ifdefined\hypertarget\hypertarget{sourcepointer-1}{}\fi
\noindent\textit{Source reference.} Original reference 0FIH: sources/stacks-weil/pinned-original-source.tex, source lines 4761--4761. Prerequisite pointer invoked at proof line(s) 5400; it names a standard fact of the same chapter that the counted argument uses, and every step of the counted proof built on it is retained verbatim in the delivered body.\par\medskip
\ifdefined\hypertarget\hypertarget{sourcepointer-2}{}\fi
\noindent\textit{Source reference.} Original reference 0FIR: sources/stacks-weil/pinned-original-source.tex, source lines 5258--5258. Prerequisite pointer invoked at proof line(s) 5391; it names a standard fact of the same chapter that the counted argument uses, and every step of the counted proof built on it is retained verbatim in the delivered body.\par\medskip
\ifdefined\hypertarget\hypertarget{sourcepointer-3}{}\fi
\noindent\textit{Source reference.} Original reference 0FIN: sources/stacks-weil/pinned-original-source.tex, source lines 5026--5026. Prerequisite pointer invoked at proof line(s) 5402; it names a standard fact of the same chapter that the counted argument uses, and every step of the counted proof built on it is retained verbatim in the delivered body.\par\medskip
\ifdefined\hypertarget\hypertarget{sourcepointer-4}{}\fi
\noindent\textit{Source reference.} Original reference 0FIS: sources/stacks-weil/pinned-original-source.tex, source lines 5336--5336. Prerequisite pointer invoked at proof line(s) 5417; it names a standard fact of the same chapter that the counted argument uses, and every step of the counted proof built on it is retained verbatim in the delivered body.\par\medskip
\begin{lemma}
\label{lemma-done}
Assume given data (D0), (D1), and (D2') satisfying axioms (A1) -- (A8).
Then the cycle class map $\gamma$ commutes with pushforward.
\end{lemma}

\begin{proof}
Let $i : Z \to X$ be as in Lemma \ref{lemma-enough-better}. Consider
the diagram
$$
\xymatrix{
E \ar[r]_j \ar[d]_\pi & X' \ar[d]^b \\
Z \ar[r]^i & X
}
$$
Let $\theta \in \CH^{r - 1}(X')$ be as in
Lemma \ref{lemma-divide-pullback-good-blowing-up}.
Then $\pi_*j^!\theta = [Z]$ in $\CH_*(Z)$ implies that
$\pi_*\gamma(j^!\theta) = 1$ by Lemma \ref{lemma-ok-for-projective-bundle}
because $\pi$ is a projective space bundle.
Hence we see that
$$
i_*(1) = i_*(\pi_*(\gamma(j^!\theta))) =
b_*j_*(j^*\gamma(\theta)) =
b_*(j_*(1) \cup \gamma(\theta))
$$
We have $j_*(1) = \gamma([E])$ by (A9). Thus this is equal to
$$
b_*(\gamma([E]) \cup \gamma(\theta)) =
b_*(\gamma([E] \cdot \theta)) =
b_*(\gamma(b^*[Z])) =
b_*b^*\gamma([Z]) = b_*(1) \cup \gamma([Z])
$$
Since $b_*(1) = 1$ by Lemma \ref{lemma-pushforward-blowup} the
proof is complete.
\end{proof}

\end{document}
